% GENERATED FILE. Edit canonical lessons, then run npm run generate:books.
% canonical-source-hash: fnv1a64:31d9cd1365ed1eaa
% canonical-lessons: PT-C414-bilheteira, PT-C414-passe, PT-C414-validar, PT-C414-revisor, PT-C414-coima

\chapter{Os bilhetes}
\label{ch:pt-os-bilhetes}
\hlchaptermodality{414}

\begin{chapteropening}
\textbf{By the end of this chapter:} \emph{I can buy and use tickets: the ticket office, the travel pass, validating it, the inspector and the fine.}

Complete the last lesson of chapter 414: I can buy and use tickets: the ticket office, the travel pass, validating it, the inspector and the fine.
\end{chapteropening}

\section[a bilheteira]{a bilheteira — the ticket office, the box office}

\label{lesson:PT-C414-bilheteira}



\begin{quote}
Before the new one: say the Portuguese for the route, then the Portuguese for to get off.
\end{quote}

\subsection*{You'll want to know: a bilheteira}
\textbf{a bilheteira} — \textquotedblleft{}the ticket office, the box office\textquotedblright{}.

\begin{quote}
A bilheteira da estação abre às seis.
\end{quote}

\begin{quote}
\emph{The station ticket office opens at six.}
\end{quote}

In Brazil: \emph{a bilheteria}. Theatres and stadiums have one too.

\begin{cousinweb}
\emph{bilheteira} is built on \emph{bilhete} 'ticket' with \emph{-eira}, a place; \emph{bilhete} came from French \emph{billet} 'note', from Old French \emph{billette} (EtymDB).
\end{cousinweb}

\begin{grammarlens}[title={What you've built}]
One more word for telling, explaining and describing at B1.
\end{grammarlens}

\subsection*{Your turn}
\begin{itemize}
\raggedright
  \item \emph{Say it:} \emph{a bilheteira}
  \item \emph{Say it:} \emph{a bilheteira}, once more
  \item \emph{Say it:} a sentence of your own with \emph{a bilheteira}
  \item \emph{Recall:} say the Portuguese for the route, then the Portuguese for to get off, then say \emph{a bilheteira} again
\end{itemize}

\begin{tcolorbox}[breakable,colback=teal!4,colframe=teal!35!black,title={Before you move on}]
What does a bilheteira mean? (\textquotedblleft{}The ticket office, the box office\textquotedblright{}.) Say it once more.
\end{tcolorbox}
\section[o passe]{o passe — the travel pass, the season ticket; the pass (in football)}

\label{lesson:PT-C414-passe}



\begin{quote}
Before the new one: say the Portuguese for to get off, then the Portuguese for the ticket office.
\end{quote}

\subsection*{You'll want to know: o passe}
\textbf{o passe} — \textquotedblleft{}the travel pass, the season ticket; the pass (in football)\textquotedblright{}.

\begin{quote}
Compro o passe todos os meses.
\end{quote}

\begin{quote}
\emph{I buy the travel pass every month.}
\end{quote}

\emph{O passe} covers buses, trams and the metro for a month.

\begin{cousinweb}
\emph{passe} is built on \emph{passar} 'to pass' (EtymDB), which is inherited from Vulgar Latin \emph{passāre}, on Latin \emph{passus} 'step'.
\end{cousinweb}

\begin{grammarlens}[title={What you've built}]
One more word for telling, explaining and describing at B1.
\end{grammarlens}

\subsection*{Your turn}
\begin{itemize}
\raggedright
  \item \emph{Say it:} \emph{o passe}
  \item \emph{Say it:} \emph{o passe}, once more
  \item \emph{Say it:} a sentence of your own with \emph{o passe}
  \item \emph{Recall:} say the Portuguese for to get off, then the Portuguese for the ticket office, then say \emph{o passe} again
\end{itemize}

\begin{tcolorbox}[breakable,colback=teal!4,colframe=teal!35!black,title={Before you move on}]
What does o passe mean? (\textquotedblleft{}The travel pass, the season ticket; the pass (in football)\textquotedblright{}.) Say it once more.
\end{tcolorbox}
\section[validar]{validar — to validate (a ticket); to approve}

\label{lesson:PT-C414-validar}



\begin{quote}
Before the new one: say the Portuguese for the ticket office, then the Portuguese for the travel pass.
\end{quote}

\subsection*{You'll want to know: validar}
\textbf{validar} — \textquotedblleft{}to validate (a ticket); to approve\textquotedblright{}.

\begin{quote}
Validamos o bilhete à entrada do metro.
\end{quote}

\begin{quote}
\emph{We validate the ticket at the metro entrance.}
\end{quote}

Machines at the gates \emph{validam} the card.

\begin{cousinweb}
\emph{validar} is a learned borrowing from Latin \emph{validāre} 'to make strong', from \emph{validus} 'strong' (EtymDB), on \emph{valēre} 'to be strong'.
\end{cousinweb}

\begin{grammarlens}[title={What you've built}]
One more word for telling, explaining and describing at B1.
\end{grammarlens}

\subsection*{Your turn}
\begin{itemize}
\raggedright
  \item \emph{Say it:} \emph{validar}
  \item \emph{Say it:} \emph{validar}, once more
  \item \emph{Say it:} a sentence of your own with \emph{validar}
  \item \emph{Recall:} say the Portuguese for the ticket office, then the Portuguese for the travel pass, then say \emph{validar} again
\end{itemize}

\begin{tcolorbox}[breakable,colback=teal!4,colframe=teal!35!black,title={Before you move on}]
What does validar mean? (\textquotedblleft{}To validate (a ticket); to approve\textquotedblright{}.) Say it once more.
\end{tcolorbox}
\section[o revisor]{o revisor — the ticket inspector}

\label{lesson:PT-C414-revisor}



\begin{quote}
Before the new one: say the Portuguese for the travel pass, then the Portuguese for to validate.
\end{quote}

\subsection*{You'll want to know: o revisor}
\textbf{o revisor} — \textquotedblleft{}the ticket inspector\textquotedblright{}.

\begin{quote}
O revisor pediu os bilhetes a toda a gente.
\end{quote}

\begin{quote}
\emph{The inspector asked everyone for their tickets.}
\end{quote}

Also a proofreader of texts.

\begin{cousinweb}
\emph{revisor} is a learned formation on Latin \emph{revīsere} 'to come back to see, to look again at', built on \emph{vidēre} 'to see' (Lewis \& Short): the person who checks the tickets again.
\end{cousinweb}

\begin{grammarlens}[title={What you've built}]
One more word for telling, explaining and describing at B1.
\end{grammarlens}

\subsection*{Your turn}
\begin{itemize}
\raggedright
  \item \emph{Say it:} \emph{o revisor}
  \item \emph{Say it:} \emph{o revisor}, once more
  \item \emph{Say it:} a sentence of your own with \emph{o revisor}
  \item \emph{Recall:} say the Portuguese for the travel pass, then the Portuguese for to validate, then say \emph{o revisor} again
\end{itemize}

\begin{tcolorbox}[breakable,colback=teal!4,colframe=teal!35!black,title={Before you move on}]
What does o revisor mean? (\textquotedblleft{}The ticket inspector\textquotedblright{}.) Say it once more.
\end{tcolorbox}
\section[a coima]{a coima — the fine (official)}

\label{lesson:PT-C414-coima}



\begin{quote}
Before the new one: say the Portuguese for to validate, then the Portuguese for the ticket inspector.
\end{quote}

\subsection*{You'll want to know: a coima}
\textbf{a coima} — \textquotedblleft{}the fine (official)\textquotedblright{}.

\begin{quote}
Quem viaja sem bilhete paga uma coima.
\end{quote}

\begin{quote}
\emph{Anyone travelling without a ticket pays a fine.}
\end{quote}

\emph{A coima} is the official word on notices; people often say \emph{a multa}.

\begin{cousinweb}
The origin of \emph{coima} is not settled. It is an old Portuguese word for a penalty, and EtymDB has no entry for it.
\end{cousinweb}

\begin{grammarlens}[title={What you've built}]
One more word for telling, explaining and describing at B1.
\end{grammarlens}

\subsection*{Your turn}
\begin{itemize}
\raggedright
  \item \emph{Say it:} \emph{a coima}
  \item \emph{Say it:} \emph{a coima}, once more
  \item \emph{Say it:} a sentence of your own with \emph{a coima}
  \item \emph{Recall:} say the Portuguese for to validate, then the Portuguese for the ticket inspector, then say \emph{a coima} again
\end{itemize}

\begin{tcolorbox}[breakable,colback=teal!4,colframe=teal!35!black,title={Before you move on}]
What does a coima mean? (\textquotedblleft{}The fine (official)\textquotedblright{}.) Say it once more.
\end{tcolorbox}
